\chapter{MARCO TEÓRICO}\label{cap:marco}
% Bloque (b), parte teorica. Ver redaccion/MAPA.md.
% Pautas de la plantilla UTEC movidas a redaccion/RUBRICA.md (P-MT).
% Fuentes: docs/03-glosario.md; fichas de docs/literatura/ (deman1999, glover1980nonlinear,
% park2015ct, selles2024marreview, deman2007catsim, wu2022xcist, ho2020denoising, song2021ddim,
% nichol2021improved, rombach2022latentdiffusion, lugmayr2022repaint, zhang2025diffboost,
% dorjsembe2024, zhang2023controlnet, kaiser2014dysmorphism, mclaren2021corridor,
% gardner2010safezones, smith2006iliosacral, zwingmann2009navigated, routt1997, wasserthal2023);
% tesis/main.tex. Fuente de cada afirmacion: redaccion/rondas/capitulo1-r00-respuesta.md.

Este capítulo reúne los conceptos que el Capítulo~\ref{cap:propuesta} usa sin volver a definirlos. Está escrito para un lector de computación que no conoce la física de la tomografía ni la cirugía de la pelvis. De cada concepto se da la definición, la notación y la fuente, y se dice en qué pieza de la propuesta interviene. El capítulo no compara trabajos entre sí, que es la función del Capítulo~\ref{cap:estado-arte}, ni describe el diseño, que es la del Capítulo~\ref{cap:propuesta}.

Las secciones van de la física de la imagen a la cirugía, al modelo generativo y a la estadística: el artefacto se explica con los conceptos de la tomografía, y el modelo 2.5D, con los del artefacto. La cirugía va antes del modelo generativo porque el muestreador, que usa sus conceptos, precede en la cadena al sintetizador (Figura~\ref{fig:pipeline}). La Sección~\ref{sec:mt-tc} define qué mide una tomografía y en qué unidades, de lo que dependen la compuerta del Objetivo~1 y la representación del sintetizador. La Sección~\ref{sec:mt-artefactos} explica cómo produce el metal su artefacto y por qué ese artefacto no queda dentro del implante. La Sección~\ref{sec:mt-iliosacra} describe la fijación iliosacra, de la que salen el corredor óseo, la escala de brecha cortical y la referencia clínica del Objetivo~2. La Sección~\ref{sec:mt-difusion} presenta los modelos de difusión con que se formula el sintetizador del Objetivo~3, y la Sección~\ref{sec:mt-estadistica}, el análisis estadístico con que se leen los resultados. La Figura~\ref{fig:mt-conceptos} relaciona cada grupo de conceptos con la pieza de la propuesta que lo usa.

\begin{figure}[htbp]
\centering
\small
\begin{tabular}{p{0.44\linewidth} c p{0.40\linewidth}}
\fbox{\parbox{0.95\linewidth}{\centering Atenuación, sinograma, reconstrucción, unidades Hounsfield, ventanas y codificación multiventana (Sección~\ref{sec:mt-tc})}} & $\rightarrow$ & \fbox{\parbox{0.93\linewidth}{\centering Compuerta de representación (Objetivo~1); canales de entrada y salida del sintetizador (Objetivo~3)}} \\ \noalign{\vspace{8pt}}
\fbox{\parbox{0.95\linewidth}{\centering Mecanismos del artefacto metálico, dominio en que se definen y métricas con que se cuantifica (Sección~\ref{sec:mt-artefactos})}} & $\rightarrow$ & \fbox{\parbox{0.93\linewidth}{\centering Banda de generación extendida $B_{\delta}$ y supuesto del dominio de imagen (Objetivo~3); métricas de apariencia (Objetivos~3 y~4); protocolo físico de comparación}} \\ \noalign{\vspace{8pt}}
\fbox{\parbox{0.95\linewidth}{\centering Tornillo iliosacro, zona segura, corredor óseo, marco de referencia y escala de brecha cortical (Sección~\ref{sec:mt-iliosacra})}} & $\rightarrow$ & \fbox{\parbox{0.93\linewidth}{\centering Muestreador de colocación (Objetivo~2) y SAP (Objetivo~4)}} \\ \noalign{\vspace{8pt}}
\fbox{\parbox{0.95\linewidth}{\centering Procesos directo e inverso, condicionamiento, difusión latente, \emph{inpainting} y modelo 2.5D (Sección~\ref{sec:mt-difusion})}} & $\rightarrow$ & \fbox{\parbox{0.93\linewidth}{\centering Autoencoder examinado en la compuerta (Objetivo~1); sintetizador (Objetivo~3)}} \\ \noalign{\vspace{8pt}}
\fbox{\parbox{0.95\linewidth}{\centering Distancia de Wasserstein-1, superioridad y equivalencia, intervalos por remuestreo (Sección~\ref{sec:mt-estadistica})}} & $\rightarrow$ & \fbox{\parbox{0.93\linewidth}{\centering Lectura de los resultados de los Objetivos~1 a~3; definición de la distancia en SAP (Objetivo~4)}} \\
\end{tabular}
\caption{Conceptos del capítulo y pieza de la propuesta que los usa.}
\label{fig:mt-conceptos}
\end{figure}

\section{Tomografía computarizada y unidades Hounsfield}\label{sec:mt-tc}

Una TC no fotografía el cuerpo: mide cuánto se atenúan los rayos X a lo largo de muchas trayectorias y reconstruye, a partir de esas medidas, un mapa de atenuación. Cada lectura del detector corresponde a un rayo de una vista, y el conjunto de lecturas ordenadas por índice forma el \textbf{sinograma} \cite{deman2007catsim}. En el modelo de adquisición de CatSim, la señal de cada índice del sinograma suma, sobre las energías del espectro, los fotones que llegarían al detector sin objeto, atenuados según la longitud que el rayo recorre en cada material. Esa atenuación depende del coeficiente de atenuación lineal de cada material a cada energía \cite{deman2007catsim}. Abadi et al.~\cite{abadi2019}, en el simulador DukeSim, y Wang et al.~\cite{wang2019cochlear}, en su simulación de implantes cocleares, calculan esa atenuación con la ley de Beer-Lambert; Wang et al. la discretizan en cinco energías.

La imagen se obtiene invirtiendo ese proceso de medida. Las simulaciones de De Man et al.~\cite{deman1999,deman2007catsim} y de Karageorgos et al.~\cite{karageorgos2024ddpm} reconstruyen con retroproyección filtrada (\emph{filtered back-projection}). Park et al.~\cite{park2015ct} señalan que ese algoritmo supone un modelo monocromático de las proyecciones, el de Radon, y que las proyecciones de un haz policromático se apartan de él de forma no lineal. Esa discrepancia reaparece en la imagen como artefacto (Sección~\ref{sec:mt-artefactos}). Este documento distingue dos dominios: el \textbf{dominio de proyección}, el del sinograma, y el \textbf{dominio de imagen}, el de los valores reconstruidos por vóxel. La cohorte de este trabajo solo contiene volúmenes ya reconstruidos (Sección~\ref{sec:datos}), de modo que lo que este trabajo mide, y lo que genera el sintetizador, ocurre en el dominio de imagen. Solo el protocolo físico de comparación, previsto en la Sección~\ref{sec:apariencia}, pasa por proyecciones simuladas.

Los valores de TC se expresan en \textbf{unidades Hounsfield} (HU), una escala normalizada de la atenuación reconstruida en la que el agua vale 0~HU y el aire $-1000$~HU por definición \cite{wu2022xcist}. La conversión entre coeficiente de atenuación y HU se apoya en el coeficiente de absorción del agua en función de la energía \cite{wang2019cochlear} \GAPLIT{fuente de referencia de física de TC que enuncie la ley de Beer-Lambert con el coeficiente de atenuación lineal, dé la definición de la escala Hounsfield como transformación de ese coeficiente y los fundamentos de la reconstrucción por retroproyección filtrada}. El protocolo de evaluación adoptado define el hueso como los vóxeles sobre 150~HU fuera de la geometría metálica \cite{peters2025hybrid}. Dos fuentes de MAR segmentan el metal clínico con 2500~HU \cite{wang2025adaptiveweighting,li2024}, umbral con que este trabajo criba la cohorte y delimita el metal clínico de sus volúmenes (Secciones~\ref{sec:datos} y~\ref{sec:sintetizador}). Los dos umbrales son definiciones operativas de las fuentes que los usan y no límites físicos del tejido. En la cohorte primaria, con el recorte por defecto, la mediana de la fracción de vóxeles a 150~HU o menos dentro de las máscaras sacras fue 0.42 a lo largo del eje del corredor más ancho. Por eso el corredor se mide sobre máscaras anatómicas y no con el umbral de 150~HU, que dejaría fuera esa fracción (Sección~\ref{sec:corredor}).

Varias de las redes de síntesis y de MAR revisadas normalizan los HU con una ventana antes de procesarlos (Sección~\ref{sec:ea-representacion}). Una \textbf{ventana} es un intervalo $[w_{\min}, w_{\max}]$ de HU que se satura y se lleva linealmente a $[0, 1]$ \cite{wang2025adaptiveweighting}:
\begin{equation}
\tilde{x} = \frac{\min\bigl(\max(x, w_{\min}), w_{\max}\bigr) - w_{\min}}{w_{\max} - w_{\min}},
\label{eq:ventana}
\end{equation}
donde $x$ es el valor de TC de un vóxel y $\tilde{x}$ su valor normalizado. Lo que cae fuera de la ventana se \textbf{satura}: todos los valores por encima de $w_{\max}$ se vuelven indistinguibles entre sí, y lo mismo ocurre bajo $w_{\min}$. Dentro de la ventana, un mismo intervalo de HU ocupa una fracción del rango normalizado inversamente proporcional al ancho de la ventana. Por eso una ventana ancha conserva más rango pero comprime los contrastes pequeños del tejido, y una estrecha los separa pero satura más valores, entre ellos los del hueso denso y el metal.

La \textbf{codificación multiventana} usa varias ventanas a la vez, una por canal, y evita así elegir entre conservar el rango y separar los contrastes. En la MAR, Wang et al.~\cite{wang2025adaptiveweighting} trabajan con tres ventanas, $[-1000, 2000]$, $[-320, 480]$ y $[-160, 240]$~HU, aunque en una cascada de etapas y no como canales de entrada (Sección~\ref{sec:ea-representacion}). El paso de ida y vuelta lleva los HU a la representación y de vuelta a HU, y su error es la diferencia, en HU, entre el valor original de cada vóxel y el recuperado.

Este trabajo usa la codificación multiventana en dos piezas: la compuerta del Objetivo~1 y el sintetizador del Objetivo~3. El Objetivo~1 examinó tres codificaciones (Sección~\ref{sec:obj1}). La primera usa como canales las ventanas de Wang et al., cuyo techo de 2000~HU queda bajo el umbral de 2500~HU del metal, de modo que el metal satura también en la ventana ancha. La segunda eleva ese techo a 20\,000~HU, y la tercera sustituye la ventana ancha por una compresión arcoseno hiperbólico sobre $[-1000, 20\,000]$~HU, que a diferencia de la Ecuación~\ref{eq:ventana} no es lineal \GAPDATO{forma y parámetros de la compresión arcoseno hiperbólico, que solo constan en el código del experimento y no en un registro de sus resultados}. La compuerta mide el error de ida y vuelta como error absoluto medio dentro del hueso (Ecuación~\ref{eq:mae}). El sintetizador usa la codificación multiventana como representación de sus canales de entrada y salida, sin la etapa del autoencoder (Sección~\ref{sec:sintetizador}), y la variante que adopta no se ha fijado \GAPDEC{qué codificación multiventana y qué precisión numérica adopta el sintetizador, pendiente de su preinscripción}.

\section{Artefactos inducidos por metal}\label{sec:mt-artefactos}

Un \textbf{artefacto metálico} no es la imagen del metal, sino el error que el metal introduce en la reconstrucción del tejido que lo rodea. Selles et al.~\cite{selles2024marreview} lo describen como rayas claras y oscuras que ocultan estructuras anatómicas. Según la misma revisión, su severidad depende del tamaño, la forma y la aleación del implante, y los materiales de osteosíntesis caen en una categoría de artefacto intermedia, asignada de forma cualitativa y sin criterio numérico.

Las fuentes revisadas atribuyen el artefacto a varios mecanismos físicos; sus listas comparten el endurecimiento del haz, la dispersión y los efectos de borde, y difieren en el resto. Selles et al. nombran como principales el endurecimiento del haz, la inanición de fotones, la dispersión y los efectos de borde. De Man et al.~\cite{deman1999} aíslan cada causa en una simulación bidimensional y concluyen que las más importantes son el endurecimiento del haz, la dispersión, el ruido y el efecto exponencial de gradiente de borde. También concluyen que todas las causas que estudian producen rayas, sin asignar a ninguna un peso numérico.

En la simulación bidimensional de De Man et al., el endurecimiento del haz y la dispersión producen rayas muy parecidas entre sí \cite{deman1999}. El \textbf{endurecimiento del haz} (\emph{beam hardening}) ocurre porque el material absorbe con preferencia los fotones de baja energía, de modo que el espectro que sale del objeto se desplaza hacia energías más altas \cite{selles2024marreview}. Park et al.~\cite{park2015ct} separan dos efectos del endurecimiento: un error de los valores dentro del propio objeto (\emph{cupping}) y rayas que corrompen la imagen fuera de la región que las produce. En esa simulación, el endurecimiento produce rayas oscuras en las direcciones de mayor atenuación y, con dos o más objetos metálicos, rayas oscuras que los conectan \cite{deman1999}. La \textbf{dispersión} (\emph{scatter}) aumenta con la alta densidad electrónica del metal \cite{selles2024marreview}. Para producir rayas apreciables, en la misma simulación basta una razón entre radiación dispersa y primaria de 0.0001, elegida de forma arbitraria \cite{deman1999}.

La \textbf{inanición de fotones} (\emph{photon starvation}) aparece cuando un rayo atraviesa tanto metal que llegan muy pocos fotones al detector, y faltan datos de proyección esenciales \cite{selles2024marreview}. De Man et al. no nombran la inanición de fotones; describen como artefacto de ruido líneas finas, alternadamente oscuras y claras, cuya intensidad depende de la atenuación total integrada a lo largo del rayo \cite{deman1999}. Según ellos, ese artefacto de ruido es no lineal, como los del endurecimiento del haz y la dispersión, y su severidad depende de cuánta dispersión y endurecimiento haya, de modo que en su simulación las causas no actúan por separado.

El \textbf{efecto exponencial de gradiente de borde} (\emph{exponential edge-gradient effect}) produce, según recuerdan De Man et al., rayas tangentes a los bordes rectos largos, y en su simulación se observan rayas que irradian desde el metal \cite{deman1999}. Selles et al. describen los efectos de borde como rayas oscuras o claras alineadas con el borde del metal. Ninguna de las dos descripciones incluye una distancia, y el artículo de De Man et al. no menciona una medida de hasta dónde llegan las rayas de ninguna de las causas que estudia.

A diferencia de los anteriores, el \textbf{volumen parcial no lineal} depende de la tercera dimensión. Glover y Pelc~\cite{glover1980nonlinear} derivan que, cuando la atenuación varía en la dirección axial dentro del espesor del haz, integrar el flujo y tomar su logaritmo vuelve no lineal la medida. En su análisis, una sola discontinuidad produce sobre todo un error local, mientras que dos o más pueden producir rayas de largo alcance que conectan estructuras. Lo estudian en hueso, sin metal y con una fuente monocromática idealizada, y De Man et al. lo dejan fuera al simular en dos dimensiones \cite{deman1999}: ninguna de las dos fuentes reúne metal y volumen parcial axial.

De estos mecanismos se siguen dos consecuencias para la síntesis. La primera es que el artefacto no queda dentro del implante: las rayas se extienden por el tejido que rodea al metal \cite{deman1999,park2015ct}. Por eso el sintetizador genera también en la banda de generación extendida $B_{\delta}$ alrededor de la máscara del implante, cuyo ancho es una convención de este trabajo porque las fuentes revisadas no dan hasta dónde llegan las rayas (Sección~\ref{sec:ea-representacion}). La segunda es que los mecanismos se definen sobre las proyecciones: De Man et al. aíslan cada causa modificando el sinograma y reconstruyendo después. Como la cohorte solo contiene volúmenes reconstruidos (Sección~\ref{sec:mt-tc}), este trabajo asume que la apariencia del artefacto puede generarse en el dominio de imagen; la Sección~\ref{sec:amenazas} discute ese supuesto.

La MAR reduce el artefacto de una TC ya adquirida, y la simulación física lo produce reproduciendo la adquisición; el sintetizador de este trabajo no hace ni lo uno ni lo otro. CatSim, el simulador del conjunto XCIST \cite{wu2022xcist} sobre el que se ejecuta el protocolo físico de Peters et al., modela el espectro policromático, el ruido cuántico y el electrónico, el volumen parcial no lineal y la dispersión \cite{deman2007catsim}. Este trabajo genera la apariencia del artefacto con un modelo aprendido, sin simular proyecciones, y prevé usar la simulación física como comparación (Sección~\ref{sec:apariencia}).

Para cuantificar el artefacto, este trabajo adopta las métricas del protocolo de Peters et al.~\cite{peters2025hybrid}, con sus nombres publicados. La \emph{streak amplitude} mide las rayas en regiones perpendiculares a ellas, como la diferencia entre el promedio del 5\,\% de desviaciones más altas respecto de la imagen sin metal del mismo caso y el del 5\,\% más bajas. La \emph{bone integrity} compara el hueso, delimitado con el umbral de 150~HU de la Sección~\ref{sec:mt-tc}, mediante el cambio de volumen y el coeficiente de Sørensen-Dice. La \emph{metal integrity} compara de forma análoga la máscara del metal, con un umbral que se adapta a cada región. Las tres se diseñaron para evaluar MAR, y su uso para evaluar síntesis exige invertirlas; la definición operativa de esa inversión está pendiente para cada métrica (Sección~\ref{sec:apariencia}). El Objetivo~3 toma la \emph{streak amplitude} como su único criterio primario.

\section{Fijación iliosacra y corredores óseos}\label{sec:mt-iliosacra}

La fijación iliosacra trata fracturas inestables del anillo pélvico posterior. Zwingmann et al.~\cite{zwingmann2009navigated} incluyen en su estudio solo fracturas de tipos~B y~C de la clasificación de Tile y Pennal. La de tipo~B es inestable para la rotación, y la de tipo~C, para la rotación y el desplazamiento vertical. En ese estudio, los tornillos se colocan por vía percutánea, y la distribución de sus grados de brecha forma la referencia clínica del Objetivo~2.

En Smith et al.~\cite{smith2006iliosacral}, el \textbf{tornillo iliosacro} entra por el ilion y llega a la primera o a la segunda vértebra sacra, de modo que atraviese tres corticales, dos del ilion y una del ala sacra. El ala sacra desciende a cada lado, hacia lateral y caudal, desde el cuerpo vertebral superior del sacro \cite{routt1997} \GAPLIT{fuente de anatomía de la pelvis que defina e ilustre la cortical ósea, el platillo vertebral, el ala sacra, el foramen neural, la tabla externa del ilion y la articulación sacroilíaca}. Un tornillo transiliosacro, en cambio, cruza las dos articulaciones sacroilíacas y sale por la tabla externa del ilion opuesto \cite{mclaren2021corridor}. Este trabajo trata como no intercambiables las medidas de corredor de los dos tipos: las de McLaren et al. son de tornillos transiliosacros, y Kaiser et al.~\cite{kaiser2014dysmorphism} eligen su umbral de 10~mm para el paso de un tornillo iliosacro. La Sección~\ref{sec:sap}, sin embargo, trata el tornillo de este trabajo como uno que entra y sale por la cortical del ilion \GAPDEC{qué tornillo representa el corredor que mide este trabajo, que el registro describe de la cortical externa de un ilion a la del ilion opuesto, y con ello si le aplican el umbral de 10~mm y la holgura radial de Kaiser et al., definidos para el tornillo iliosacro, y si es del mismo tipo que los tornillos de la referencia clínica, que Zwingmann et al. llaman iliosacros y transiliosacros}.

La colocación percutánea se guía con imagen intraoperatoria. Routt et al.~\cite{routt1997} describen la técnica con el paciente en decúbito supino y fluoroscopia en tres planos. Zwingmann et al. comparan la guía fluoroscópica, que llaman convencional, con navegación computarizada sobre una adquisición tridimensional intraoperatoria hecha con un equipo que rota $190^{\circ}$ alrededor del campo quirúrgico. De esa comparación salen la serie navegada y la serie convencional. En ambas, la posición final de cada tornillo se evaluó sobre una TC posoperatoria.

El tornillo debe quedar dentro del hueso en una región del sacro superior, la \textbf{zona segura}. Routt et al. la describen en el ala sacra, limitada por delante y por arriba por la cortical del ala, y por detrás y por abajo por el túnel del foramen neural. Cerca de ella quedan la raíz nerviosa de L5 sobre el ala, el canal espinal y los vasos ilíacos por delante del ala, y un tornillo que sale de la zona se acerca a esas estructuras \cite{routt1997}. Gardner et al.~\cite{gardner2010safezones} representan la zona como dos conos unidos por la punta y miden su sección transversal mínima donde los conos se encuentran, en un plano ortogonal al eje de la zona. Este trabajo no usa la zona segura como variable: la presenta como antecedente del corredor óseo, que es lo que el muestreador mide (Sección~\ref{sec:muestreador}).

El \textbf{corredor óseo} extiende esa idea a toda la trayectoria: es la región de hueso por la que puede pasar un tornillo sin atravesar la cortical. McLaren et al.~\cite{mclaren2021corridor} lo resumen en un solo número por segmento sacro. Colocan una recta dentro de los contornos corticales y aumentan su diámetro hasta que toca y atraviesa la cortical en al menos tres puntos, lo que da el diámetro máximo de la trayectoria. Gardner et al., Kaiser et al.~\cite{kaiser2014dysmorphism} y McLaren et al. emplean un umbral de 10~mm sobre el tamaño del corredor. Los tres lo toman de trabajos anteriores, y McLaren et al. declaran además que el corredor mínimo no está establecido. Kaiser et al. lo justifican como una holgura de 1 a 2~mm alrededor de un tornillo de 6.3 a 8~mm, que la Sección~\ref{sec:corredor} lee como holgura radial por lado. Esa sección describe también cómo este trabajo mide el corredor, sobre máscaras anatómicas de TotalSegmentator \cite{wasserthal2023}, y qué criterio de viabilidad adopta en lugar de ese umbral.

Expresar la pose de un tornillo exige un marco de referencia anatómico que se pueda calcular sobre la TC. Kaiser et al. reformatean el volumen a lo largo del eje del sacro, perpendicular al platillo superior de S1. Miden la angulación del corredor contra la línea que une las crestas ilíacas y contra la que une las espinas ilíacas posteriores. El marco incluye además una regla de longitud útil, que exige al menos 5~mm hasta la cortical a cada lado del eje y que la Sección~\ref{sec:corredor} lee como margen cortical. Este trabajo expresa sus poses en ese marco, y la Sección~\ref{sec:marco-metal} examina si el platillo de S1, las crestas y las espinas ilíacas se pueden localizar en volúmenes con metal.

Smith et al.~\cite{smith2006iliosacral} juzgan la posición de un tornillo frente a una posición ideal. La definen enteramente dentro de los márgenes corticales del sacro, paralela al platillo correspondiente y a media altura entre el platillo superior de S1 y el foramen neural de S1. Distinguen tres tipos de perforación según la cortical que el tornillo atraviesa: la anterior del sacro, la del canal espinal y la de los platillos vertebrales. Smith et al. gradúan esa perforación, que en este documento se llama brecha cortical, en cuatro niveles. El grado~0 es la ausencia de perforación, el grado~1 una perforación menor que 2~mm, el grado~2 una de 2 a 4~mm y el grado~3 una mayor que 4~mm. En este trabajo, el grado se calcula a partir de la protrusión del implante fuera de una envolvente ósea segmentada, que aproxima esa perforación (Secciones~\ref{sec:sap} y~\ref{sec:amenazas}).

Las fuentes revisadas no juzgan la colocación con una misma escala. Smith et al. declaran que toman su escala de la clasificación de tornillos pediculares, y le suman una escala angular \GAPDEC{si SAP incorpora la dimensión angular de la escala de Smith et al. o declara que usa solo la de perforación, y con qué justificación}. Zwingmann et al. aplican la escala de perforación (Sección~\ref{sec:sap}). Hinsche et al.~\cite{hinsche2002fluoroscopy} usan, en cambio, una definición binaria: consideran insegura una colocación cuya trayectoria perfora una de las corticales comprometiendo estructuras neurovasculares, sin grados de profundidad.

La escala es ordinal: sus grados están ordenados por la profundidad de la perforación, pero no cubren intervalos del mismo ancho. Los grados~1 y~2 cubren 2~mm cada uno, y el grado~3 no tiene límite superior. Tratar los grados como equiespaciados, como hace la distancia de la Sección~\ref{sec:mt-estadistica}, es entonces una convención y no una propiedad de la escala (Sección~\ref{sec:sap}).

La escala de brecha y la viabilidad del corredor son dos de los tres componentes de SAP; el tercero es la fracción por zona de densidad. La motivación de esa fracción es la heterogeneidad de la masa ósea pélvica que describen Arand et al.~\cite{arand2019pelvicring} con un modelo medio de los valores de gris del anillo pélvico. En ese modelo, los valores del ala sacra son más bajos que los del cuerpo de S1 (Sección~\ref{sec:poses}). SAP reporta esa fracción de forma descriptiva, sin compararla con la referencia clínica, y su definición operativa está pendiente \GAPDEC{fuente y definición operativa de la fracción por zona de densidad de SAP}.

\section{Modelos de difusión y síntesis condicionada}\label{sec:mt-difusion}

Un modelo generativo aprende, a partir de ejemplos, una distribución de la que extrae muestras nuevas. Los modelos de difusión son una familia de ellos, y Kazerouni et al.~\cite{kazerouni2023diffusionsurvey} revisan su uso en imagen médica. Ho et al.~\cite{ho2020denoising} formulan la generación como una cadena de Markov que aprende a invertir un \textbf{proceso directo} que añade ruido gaussiano en $n_{\max}$ pasos; en sus experimentos, $n_{\max} = 1000$. Dorjsembe et al.~\cite{dorjsembe2024} resumen el propósito de ese proceso directo: borrar por completo la estructura original de los datos. La imagen ruidosa del paso $n$ se obtiene de la original en una sola operación, como la escriben Zhang et al.~\cite{zhang2025diffboost} en su método DiffBoost y Dorjsembe et al.:
\begin{equation}
\mathbf{x}_n = \sqrt{\bar{\gamma}_n}\,\mathbf{x}_0 + \sqrt{1 - \bar{\gamma}_n}\,\boldsymbol{\xi}, \qquad \boldsymbol{\xi} \sim \mathcal{N}(\mathbf{0}, \mathbf{I}),
\label{eq:difusion-directa}
\end{equation}
donde $\mathbf{x}_0$ es la imagen original, $\boldsymbol{\xi}$ un ruido gaussiano estándar del mismo tamaño y $\bar{\gamma}_n$ un coeficiente que fija el \textbf{calendario de ruido}. Las fuentes escriben ese coeficiente como $\bar{\alpha}_t$; aquí se cambia de símbolo porque $\alpha$ y $t$ designan otras magnitudes en el Capítulo~\ref{cap:propuesta}. Nichol y Dhariwal~\cite{nichol2021improved} definen un calendario coseno para ese coeficiente y obtienen, de sus cocientes sucesivos, las varianzas $\beta_n$ de cada paso; Zhang et al.~\cite{zhang2025diffboost} las sitúan en el intervalo $(0, 1)$.

En la formulación de Ho et al., la red no predice la imagen limpia, sino el ruido añadido, con un objetivo simplificado que Rombach et al.~\cite{rombach2022latentdiffusion} escriben como
\begin{equation}
\mathcal{L} = \mathbb{E}_{\mathbf{x}_0,\, \boldsymbol{\xi},\, n}\Bigl[\bigl\lVert \boldsymbol{\xi} - \boldsymbol{\xi}_{\theta}(\mathbf{x}_n, n) \bigr\rVert_2^2\Bigr],
\label{eq:perdida-difusion}
\end{equation}
con $n$ tomado de forma uniforme entre 1 y $n_{\max}$, donde $\boldsymbol{\xi}_{\theta}$ es la red con parámetros $\theta$. Para generar, el \textbf{proceso inverso} parte de ruido puro y aplica la red paso a paso, del último al primero. En Ho et al. y en Song et al.~\cite{song2021ddim}, esa red es una U-Net. Ronneberger et al.~\cite{ronnenberger2015unet} la definen como una red completamente convolucional con una ruta que contrae la resolución y otra simétrica que la expande, y que combina las características de alta resolución de la primera con la segunda. El borrador del diseño del sintetizador propone también una U-Net, arquitectura que aún no se ha fijado (Sección~\ref{sec:sintetizador}).

El costo de generar crece con el número de pasos del proceso inverso. Song et al. definen una familia de procesos no markovianos que conserva el objetivo de entrenamiento de Ho et al. y permite muestrear de forma determinista y con menos pasos, sin reentrenar la red. Reportan una generación entre 10 y 50 veces más rápida, en tiempo de reloj, que la del muestreo de Ho et al. Con ello, en un modelo entrenado con el objetivo de Ho et al., el \textbf{procedimiento de muestreo} se elige por separado del entrenamiento. Las tres fuentes evalúan imágenes naturales y ninguna trabaja con TC ni con metal \cite{ho2020denoising,song2021ddim,nichol2021improved}, de modo que la aceleración que reportan Song et al. no está medida sobre TC. El borrador del diseño del sintetizador propone el muestreo determinista de Song et al. y no menciona el calendario de ruido; ambos quedan pendientes, con el resto de su entrenamiento y muestreo (Sección~\ref{sec:sintetizador}).

Un modelo condicionado genera a partir de una condición, además del ruido. Rombach et al. distinguen dos vías: concatenar a la entrada de la red una información espacialmente alineada con la imagen, como una máscara, o introducir condiciones no espaciales, como un texto, mediante atención cruzada. Dorjsembe et al. siguen la primera vía y concatenan la máscara como canal adicional en cada paso. Zhang et al.~\cite{zhang2023controlnet} proponen ControlNet, que congela un modelo de difusión preentrenado y entrena una copia de su codificador que recibe la condición espacial. Esa copia se une al modelo congelado mediante convoluciones inicializadas en cero.

La difusión latente separa la generación en dos etapas \cite{rombach2022latentdiffusion}: un autoencoder comprime la imagen a un espacio latente de menor resolución, y un modelo de difusión opera en ese espacio. Al generar, el decodificador del autoencoder devuelve la imagen en una sola pasada. Rombach et al. entrenan el autoencoder con una pérdida perceptual y una adversarial, para evitar el desenfoque que producen las pérdidas por píxel, y describen su compresión como una etapa que elimina el detalle de alta frecuencia. También advierten que la capacidad de reconstrucción del autoencoder puede volverse un cuello de botella cuando la aplicación exige exactitud fina por píxel. Este trabajo lee como una aplicación de ese tipo conservar los HU del hueso con un error absoluto medio menor que 25~HU (Sección~\ref{sec:obj1}). Por eso el Objetivo~1 midió la ida y vuelta por el autoencoder como condición para adoptar la difusión latente, que el veredicto negativo descartó. Como ControlNet presupone un modelo base congelado y su espacio latente, quedó descartado con ella (Sección~\ref{sec:sintetizador}).

El \emph{inpainting} regenera una región enmascarada de la imagen y conserva el resto. En las fuentes revisadas, el \emph{inpainting} con un modelo de difusión sigue una de dos formas. Lugmayr et al.~\cite{lugmayr2022repaint} usan un modelo preentrenado sin condición y alteran solo el muestreo. En cada paso combinan la región conocida, muestreada de la entrada, con la desconocida, generada por el modelo, y además retroceden y avanzan en el tiempo de difusión para armonizar el borde entre ambas. La otra forma es entrenar un modelo condicionado, como el de \emph{inpainting} que Rombach et al. condicionan por concatenación. El sintetizador de este trabajo sigue esta segunda forma: recibe el parche con la región de generación $G$ borrada y las máscaras, genera solo dentro de $G$ y copia el resto del volumen de origen (Sección~\ref{sec:sintetizador}). No se ha fijado si su procedimiento de muestreo toma algo del de Lugmayr et al. o del método LeFusion de Zhang et al.~\cite{zhang2025lefusion} \GAPDEC{descripción y atribución del procedimiento de muestreo de la difusión del sintetizador frente a los de Lugmayr et al. y de LeFusion, pendiente en el registro del proyecto}.

El sintetizador es \textbf{2.5D}: opera sobre cortes axiales, pero recibe como contexto cortes contiguos al que genera, sin ser un modelo tridimensional completo. El borrador de su diseño propone tres cortes contiguos y generar el central \GAPDEC{número de cortes contiguos de la entrada 2.5D del sintetizador, que solo propone el borrador de su diseño, pendiente de preinscripción}. Como el volumen parcial no lineal nace de integrar la atenuación dentro del espesor del haz (Sección~\ref{sec:mt-artefactos}), este trabajo lee que unos pocos cortes vecinos aportan información axial pero no reproducen esa integración. El sintetizador, que opera sobre valores reconstruidos, no contiene entonces ese mecanismo y a lo sumo puede aprender la apariencia que deja en los cortes.

\section{Marco estadístico de la evaluación}\label{sec:mt-estadistica}

Los resultados de este trabajo son de tres tipos, y cada uno se analiza de forma distinta. El Objetivo~1 compara un error medio con un umbral fijado de antemano, el Objetivo~2 compara dos distribuciones de grados ordinales y el Objetivo~3 compara brazos pareados sobre los mismos pacientes. En el Objetivo~1, la unidad de análisis es el paciente: el error se promedia por paciente. En el Objetivo~2, la distribución reúne todas las poses de la cohorte, y la referencia clínica cuenta tornillos y no pacientes, de modo que ninguna de las dos distribuciones tiene al paciente como unidad (Sección~\ref{sec:amenazas}). En el Objetivo~3, las pruebas se parean por paciente, y la forma de agregar antes las poses y las regiones de medición de rayas está pendiente \GAPDEC{cómo se agregan por paciente las poses y las regiones de medición de rayas antes de las pruebas}.

Para comparar distribuciones de grados, este trabajo usa la \textbf{distancia de Wasserstein-1}. Sobre los cuatro grados de brecha, tratados como equiespaciados, es la suma, de un grado a otro, de las diferencias absolutas entre las dos funciones de distribución acumulada (Ecuación~\ref{eq:w1}). Su unidad es el grado. Si una distribución pone toda su masa un grado más arriba que la otra, la distancia vale 1; si una está entera en el grado~0 y la otra en el grado~3, vale 3, el máximo con cuatro grados. A diferencia de comparar solo la fracción de grado~0, la distancia usa la distribución completa, y dos distribuciones con la misma fracción de grado~0 difieren si reparten distinto el resto \GAPLIT{fuente de referencia de la distancia de Wasserstein-1 y de su forma para distribuciones sobre una recta}.

El Objetivo~3 plantea dos preguntas sobre sus brazos pareados. Una prueba de \textbf{superioridad} pregunta si una condición supera a la otra, y una de \textbf{equivalencia}, si la diferencia entre ambas es tan pequeña que se pueden tratar como intercambiables. Una diferencia no significativa no responde la segunda pregunta, porque puede deberse a una muestra pequeña y no a una diferencia pequeña. Por eso la equivalencia solo se concluye cuando el intervalo de confianza del 90\,\% de la diferencia pareada cae entero dentro de un margen $[-\Delta, +\Delta]$ fijado de antemano.

La prueba de superioridad del Objetivo~3 es la de rangos con signo de Wilcoxon, pareada y de una cola. Para la comparabilidad con el protocolo físico, el diseño prevé dos pruebas unilaterales de equivalencia, con un margen definido por la variabilidad entre semillas del sintetizador y entre corridas repetidas del protocolo físico, y aún sin preinscribir (Sección~\ref{sec:apariencia}) \GAPDEC{si la comparación de equivalencia frente al protocolo físico se mantiene o pasa a trabajo futuro, contingencia de plazo registrada por la autora}. Con muestras pequeñas, esa comparación puede no concluir ni equivalencia ni diferencia.

La compuerta del Objetivo~1 acompaña su error con un intervalo de confianza por remuestreo de pacientes. El error que decide es la media, sobre los pacientes, del error de cada paciente. Esa media se vuelve a calcular sobre muestras de pacientes tomadas con reemplazo, y el intervalo del 95\,\% sale de la distribución de esas medias (Sección~\ref{sec:obj1}). Se remuestrean pacientes porque el paciente es la unidad de análisis del Objetivo~1 \GAPLIT{fuentes de referencia de la prueba de rangos con signo de Wilcoxon, de las dos pruebas unilaterales de equivalencia y de su relación con el intervalo de confianza del 90\,\%, y del intervalo de confianza por remuestreo}.

La protección contra ajustar el resultado a los datos depende de lo que se fija antes de verlos y de cuántas oportunidades de aprobar da la regla. La regla de la compuerta se fijó antes de correr la prueba que decide, y la distribución de poses del muestreador y su métrica, antes de calcular cualquier distancia (Secciones~\ref{sec:obj1} y~\ref{sec:poses}). En ambos casos hubo, sin embargo, decisiones posteriores a ver datos: esa regla se fijó conociendo un resultado exploratorio desfavorable que incluía a los pacientes de prueba, y en el Objetivo~2 hubo tres (Sección~\ref{sec:amenazas}). La regla de la compuerta aprueba si lo hace cualquiera de sus seis combinaciones de autoencoder y codificación; esa \textbf{multiplicidad} sin corregir solo puede favorecer un veredicto aprobatorio y, como el veredicto fue negativo, no lo compromete. La compuerta se extendió después al espacio latente del generador de volúmenes de TC de Guo et al.~\cite{guo2025maisi}, con la misma regla y los mismos pacientes de prueba, lo que eleva a siete las combinaciones que pueden aprobar. El veredicto de esa extensión se presenta en el capítulo de resultados y, según la decisión que la autorizó, no cambia el diseño del Objetivo~3. La multiplicidad de siete combinaciones solo podría comprometer ese veredicto si fuera aprobatorio (Sección~\ref{sec:amenazas}).

El error absoluto medio y la raíz del error cuadrático medio, los dos errores por vóxel del documento, no son intercambiables. El error absoluto medio promedia el valor absoluto de las diferencias en HU (Ecuación~\ref{eq:mae}). La raíz del error cuadrático medio promedia sus cuadrados y toma la raíz, de modo que, para los mismos residuos, nunca es menor que el error absoluto medio. De ahí que los errores que reporta la literatura de MAR, expresados como raíz del error cuadrático medio, sirvan al Objetivo~1 como orden de magnitud y no como equivalente de su error absoluto medio (Sección~\ref{sec:ea-representacion}).
